\section{方法}
\label{sec:method}
\label{sec:formulation}

\subsection{总体流程与问题定义}

层边界的原子位置决定当前门的执行条件及后续重排的起点。GA-LK以相邻双比特层之间的状态转换为优化单位，联合确定目标门位及方向、跨层驻留和回迁存储位。记$s_\ell$为目标层$\mathcal G_\ell$的输入状态，包含映射$\pi_\ell$、纠缠区驻留原子、AOD资源、逐原子时钟和门依赖。$\mathcal G_\ell$中的双比特门互不共享原子，完成相应重排与门执行后得到下一层的输入状态$s_{\ell+1}$。

GA-LK在动态放置阶段求解层边界联合优化，门调度、初始布局、原子路由和分区架构中间表示（zoned architecture intermediate representation, ZAIR）代码生成沿用分区编译流程\cite{lin2025zac}。调度将输入划分为交替的单、双比特门层，初始布局确定原子到阱位的映射。层边界优化比较各候选的轨迹分批与物理代价，代码生成阶段将选定方案按执行顺序转换为ZAIR指令。各阶段及其层间状态关系见图~\ref{fig:overall-framework}。

\begin{figure*}[t]
  \centering
  \includegraphics[width=\textwidth]{../build/paper_zh/figures/overall_framework.pdf}
  \caption{分区中性原子电路的端到端编译流程。(a) 输入电路与门执行区域，(b) 分层调度，(c) 初始映射$\pi_0$；(d) 层边界联合决策产生候选映射$\pi'_\ell$，经物理状态转换后从$s_{\ell+1}$展开有限视界评价；(e) 回迁与入区轨迹按兼容性形成重排批次$\mathcal B_0$--$\mathcal B_2$；(f) 依次输出ZAIR指令并独立验证。$\mathcal E$与$\mathcal S$分别表示纠缠区和存储区。}
  \label{fig:overall-framework}
\end{figure*}

令$\mathcal Q$为逻辑原子集合，$\mathcal S$和$\mathcal E$分别为存储区与纠缠区的单原子阱位，$\mathcal V=\mathcal S\cup\mathcal E$。纠缠区门位对构成集合$\mathcal P^{\rm gate}\subseteq\mathcal E\times\mathcal E$。边界映射的单射约束要求每个阱位至多容纳一个原子：
\begin{equation}
\pi_\ell:\mathcal Q\rightarrow\mathcal V,\qquad
\pi_\ell(q)\ne\pi_\ell(q'),~q\ne q'.
\label{eq:mapping}
\end{equation}
对目标门$g_{\ell,j}=(u_j,v_j)\in\mathcal G_\ell$，集合$\Omega_{\ell,j}\subseteq\mathcal P^{\rm gate}$给出合法门位及方向。候选后映射$\pi'_{\ell}$还需满足
\begin{equation}
\begin{aligned}
(\pi'_{\ell}(u_j),\pi'_{\ell}(v_j))&\in\Omega_{\ell,j},
&&\forall g_{\ell,j}\in\mathcal G_\ell,\\
|\pi_{\ell}^{\prime-1}(\mathcal E)|&\le C_{\mathcal E},
&|\pi_{\ell}^{\prime-1}(\mathcal S)|&\le C_{\mathcal S},
\end{aligned}
\label{eq:placement-constraints}
\end{equation}
其中$C_{\mathcal E}$和$C_{\mathcal S}$为区域容量。集合$\mathcal D_\ell$包含当前边界需确定去向的纠缠区原子，包括暂不参与目标层及后续不再参与门的原子。对于仅含一个双比特门的目标层，已在纠缠区的参与原子也可纳入$\mathcal D_\ell$，比较区内调整与经存储区中转两种路径。Raman单比特旋转按门依赖计入逐原子时钟\cite{lin2025zac}。

\subsection{联合表示与物理可行域}

目标门位与驻留原子共享纠缠区空间，两类决策须共同满足占位约束。层边界方案将门位与驻留选择编码为
\begin{equation}
\mathbf x_\ell=(z_{\ell,1},\ldots,z_{\ell,m},b_{\ell,1},\ldots,b_{\ell,k}),
\label{eq:chromosome}
\end{equation}
其中$m=|\mathcal G_\ell|$、$k=|\mathcal D_\ell|$；$z_{\ell,j}\in\{0,\ldots,|\Omega_{\ell,j}|-1\}$选择门位及方向，$b_{\ell,i}=0$与$b_{\ell,i}=1$分别表示驻留与回迁。两个编码段确定目标门位和待回迁集合，补充存储位分配后得到候选后状态。图~\ref{fig:joint-ga}以两组目标门和$q_3$的回迁说明编码与轨迹的对应关系。

\begin{figure}[!t]
  \centering
  \resizebox{\columnwidth}{!}{% Shared vector-figure vocabulary for the Chinese IEEE manuscript.
% Every figure in this directory inputs this file, so the manuscript only
% needs the standard TikZ package and the libraries already loaded by
% paper_zh.tex.  Keep the guard: figures may be input more than once while
% editing or when a stand-alone smoke-test document is built.
\ifdefined\GALKFigureStyleLoaded\else
\def\GALKFigureStyleLoaded{1}

% Baseline-derived visual tokens: ICCAD uses a restrained blue/orange device
% palette, while ZAC uses conventional black axes and a small fixed plot set.
\definecolor{galkInk}{HTML}{231F20}
\definecolor{galkMuted}{HTML}{666666}
\definecolor{galkGrid}{HTML}{D6D6D6}
\definecolor{galkPaper}{HTML}{FFFFFF}
% Semantic colors are fixed across every figure: orange=entanglement zone,
% blue=storage zone, black=atom identity, red=invalidity, and blue dashes=future
% prediction.  Plot and search-group colours are deliberately independent of
% the two physical-zone colours.
\definecolor{galkStorage}{HTML}{0064BC}
\definecolor{galkStorageFill}{HTML}{E8F1F7}
\definecolor{galkEntangle}{HTML}{E37323}
\definecolor{galkEntangleFill}{HTML}{F9ECE3}
\definecolor{galkResident}{HTML}{485CB3}
\definecolor{galkResidentFill}{HTML}{EEF0F8}
\definecolor{galkGood}{HTML}{A4AF07}
\definecolor{galkBad}{HTML}{C84E45}
\definecolor{galkLook}{HTML}{0064BC}
\definecolor{galkData}{HTML}{315E8A}
\definecolor{galkTail}{HTML}{B44B27}
\definecolor{galkGroupOne}{HTML}{4C78A8}
\definecolor{galkGroupTwo}{HTML}{7A7A7A}
\definecolor{galkGroupThree}{HTML}{9C755F}

\tikzset{
  % Shared engineering-figure contract.  Physical snapshots must place an
  % orange entanglement band above a blue storage field; only their scale may
  % change.  Labels are aligned at the upper-left of both regions.
  galk/base/.style={
    font=\rmfamily\fontsize{9.0pt}{10.0pt}\selectfont,
    text=galkInk,
    line cap=butt,
    line join=miter,
    >=Stealth
  },
  galk/panel title/.style={
    font=\rmfamily\bfseries\fontsize{9.2pt}{10.2pt}\selectfont,
    text=galkInk,
    anchor=west
  },
  galk/label/.style={
    font=\rmfamily\fontsize{9.0pt}{10.0pt}\selectfont,
    text=galkInk
  },
  galk/small label/.style={
    font=\rmfamily\fontsize{8.7pt}{9.7pt}\selectfont,
    text=galkInk
  },
  galk/zone label/.style={
    font=\rmfamily\bfseries\fontsize{8.7pt}{9.7pt}\selectfont,
    text=galkInk,
    anchor=north west,
    inner sep=.5pt
  },
  galk/entanglement zone/.style={
    draw=galkEntangle,
    fill=galkEntangleFill,
    line width=.40pt
  },
  galk/storage zone/.style={
    draw=galkStorage,
    fill=galkStorageFill,
    line width=.40pt
  },
  galk/site/.style={
    circle,
    draw=galkInk!68,
    fill=white,
    line width=.36pt,
    minimum size=1.35mm,
    inner sep=0pt
  },
  galk/occupied atom/.style={
    circle,
    draw=galkInk,
    fill=galkInk,
    line width=.36pt,
    minimum size=1.55mm,
    inner sep=0pt
  },
  galk/target site/.style={
    circle,
    draw=galkInk!60,
    fill=galkInk!10,
    line width=.34pt,
    minimum size=1.55mm,
    inner sep=0pt
  },
  galk/q label/.style={
    font=\rmfamily\itshape\fontsize{8.5pt}{9.5pt}\selectfont,
    text=galkInk,
    inner sep=.25pt
  },
  galk/rydberg pair/.style={
    draw=galkInk!62,
    densely dashed,
    line width=.36pt
  },
  galk/accepted path/.style={
    -{Stealth[length=.95mm,width=.62mm]},
    draw=galkInk,
    line width=.48pt
  },
  galk/alternative path/.style={
    -{Stealth[length=.95mm,width=.62mm]},
    draw=galkMuted,
    densely dashed,
    line width=.48pt
  },
  galk/future path/.style={
    -{Stealth[length=.95mm,width=.62mm]},
    draw=galkLook,
    densely dashed,
    line width=.48pt
  },
  galk/invalid path/.style={
    -{Stealth[length=.95mm,width=.62mm]},
    draw=galkBad,
    densely dashed,
    line width=.50pt
  }
}
\fi

\begin{tikzpicture}[
  galk/base,x=1.04cm,y=1cm,
  panel title/.style={galk/panel title,font=\rmfamily\bfseries\fontsize{9.6pt}{10.6pt}\selectfont},
  label/.style={font=\rmfamily\fontsize{9.2pt}{10.2pt}\selectfont},
  note/.style={font=\rmfamily\itshape\fontsize{9.0pt}{10.0pt}\selectfont,
               text=galkMuted},
  tiny/.style={font=\rmfamily\fontsize{9.0pt}{10.0pt}\selectfont,
               text=galkInk},
  rule/.style={draw=galkInk!30,line width=.38pt},
  flow/.style={-{Stealth[length=1.0mm,width=.66mm]},draw=galkInk,line width=.50pt},
  cell/.style={draw=galkInk!72,fill=white,line width=.38pt,
               minimum width=.77cm,minimum height=.43cm,inner sep=.5pt},
  gate cell/.style={cell,draw=galkEntangle,fill=white},
  bit cell/.style={cell,draw=galkResident,fill=galkResidentFill},
  trap/.style={galk/site},
  atom/.style={galk/occupied atom,minimum size=1.85mm},
  target/.style={galk/target site,minimum size=1.85mm},
  q label/.style={galk/q label,font=\rmfamily\itshape\fontsize{9.0pt}{10.0pt}\selectfont}
]
  \path[use as bounding box] (0,.40) rectangle (8.35,7.35);
  \draw[rule] (.10,5.12)--(8.25,5.12);
  \draw[rule] (.10,4.05)--(8.25,4.05);

  % ---------------------------------------------------------------------
  % A joint chromosome and the induced injective storage assignment.
  \node[panel title] at (.12,7.15) {(a) Encoding a placement candidate};
  \node[note] at (1.93,6.83) {chromosome $\mathbf x_\ell$};
  \foreach \i/\head/\value/\cellstyle in {
      0/{$z_{\ell,1}$}/{$T_1^{+}$}/gate cell,
      1/{$z_{\ell,2}$}/{$T_2^{+}$}/gate cell,
      2/{$b_{\ell,1}$}/{1}/bit cell} {
    \node[tiny] at (.76+1.17*\i,6.46) {\head};
    \node[\cellstyle] at (.76+1.17*\i,6.10) {\value};
  }
  \node[tiny] at (.76,5.68) {$(q_0,q_2)$};
  \node[tiny] at (1.93,5.68) {$(q_1,q_4)$};
  \node[tiny] at (3.10,5.68) {$q_3$};
  \node[note] at (1.93,5.30) {gate atoms: left to right};

  \draw[flow] (3.58,6.10)--(4.12,6.10);
  \node[label] at (6.14,6.90) {$K$-best storage assignment};
  \node[tiny,anchor=west] at (4.25,6.53)
    {$q_3$ to storage};
  \node[atom] at (4.47,6.12) {};
  \node[q label,anchor=north] at (4.47,5.99) {$q_3$};
  \foreach \y/\site in {6.55/{$s_1$},6.10/{$s_2$},5.65/{$s_3$}} {
    \node[trap,minimum size=2.05mm] at (7.39,\y) {};
    \node[tiny,anchor=west] at (7.56,\y) {\site};
  }
  \draw[galkInk!34,line width=.34pt] (4.57,6.12)--(7.29,6.55);
  \draw[galkInk!34,line width=.34pt] (4.57,6.12)--(7.29,6.10);
  \draw[galkInk!34,line width=.34pt] (4.57,6.12)--(7.29,5.65);
  \draw[galkInk,line width=.62pt] (4.57,6.12)--(7.29,6.55);
  \node[label,anchor=east] at (8.15,5.30)
    {joint candidate $(\mathbf x_\ell,a_\ell)$};

  % ---------------------------------------------------------------------
  % Enumeration and genetic search share the same executable evaluation.
  \node[panel title] at (.12,4.88) {(b) Search with a common objective};
  \node[tiny] at (2.20,4.39) {Enumeration / Genetic search};
  \draw[flow] (4.46,4.39)--(4.93,4.39);
  \node[tiny] at (6.67,4.39) {Candidate evaluation};

  % ---------------------------------------------------------------------
  % Release the target sites, including temporary storage for the q1/q2 swap.
  \def\routeInitial{0/.36/1.91,1/.58/1.91,2/1.02/1.91,3/1.24/1.91,4/.63/.70}
  \def\routeStaged{0/.36/1.91,1/.28/1.08,2/1.02/1.91,3/1.33/1.08,4/.63/.70}
  \def\routeFinal{0/.36/1.91,1/1.02/1.91,2/.58/1.91,3/1.33/1.08,4/1.24/1.91}
  \def\routeMoves{0/1/.58/1.91/.28/1.08,0/3/1.24/1.91/1.33/1.08,1/2/1.02/1.91/.58/1.91,2/1/.28/1.08/1.02/1.91,3/4/.63/.70/1.24/1.91}
  \begin{scope}[shift={(0,.30)},x=1.18cm]
  \node[panel title] at (.12,3.40) {(c) Decoding and ordered execution};

  \foreach \left in {.18,2.93,5.68} {
    \path[galk/entanglement zone]
      (\left,1.63) rectangle (\left+1.60,2.80);
    \path[galk/storage zone]
      (\left,.20) rectangle (\left+1.60,1.55);
    \node[tiny] at (\left+.47,2.55) {$T_1$};
    \node[tiny] at (\left+1.13,2.55) {$T_2$};
    \draw[galk/rydberg pair] (\left+.47,1.91) ellipse (.20 and .13);
    \draw[galk/rydberg pair] (\left+1.13,1.91) ellipse (.20 and .13);
    \foreach \dx in {.36,.58,1.02,1.24}
      \node[trap] at (\left+\dx,1.91) {};
    \foreach \dx in {.28,.63,.98,1.33}
      \foreach \yy in {.70,1.08}
        \node[trap] at (\left+\dx,\yy) {};
  }
  \node[label] at (.98,3.00) {$\mathbf{s}_\ell$};
  \node[label] at (3.73,3.00) {After staging};
  \node[label] at (6.48,3.00) {$\mathbf{s}_{\ell+1}$};

  % Source state before G_l placement: two occupied pairs and q4 in storage.
  \foreach \qq/\xx/\yy in \routeInitial
    \node[atom] at (.18+\xx,\yy) {};
  \node[q label,font=\fontsize{8.5pt}{9.5pt}\selectfont] at (.65,2.20) {$q_0,q_1$};
  \node[q label,font=\fontsize{8.5pt}{9.5pt}\selectfont] at (1.31,2.20) {$q_2,q_3$};
  \draw[galkInk,line width=.46pt] (.54,1.91)--(.76,1.91);
  \draw[galkInk,line width=.46pt] (1.20,1.91)--(1.42,1.91);
  \node[q label,anchor=north] at (.81,.58) {$q_4$};

  % Intermediate state: q1 is temporarily stored; q3 remains in storage.
  \foreach \qq/\xx/\yy in \routeStaged
    \node[atom] at (2.93+\xx,\yy) {};
  \node[q label] at (3.40,2.20) {$q_0$};
  \node[q label] at (4.06,2.20) {$q_2$};
  \node[q label,anchor=north] at (3.56,.58) {$q_4$};
  \node[q label,anchor=south] at (3.21,1.20) {$q_1$};
  \node[q label,anchor=south east] at (4.44,1.20) {$q_3@s_1$};

  % Target state: q0-q2 and q1-q4 occupy the selected gate pairs.
  \foreach \qq/\xx/\yy in \routeFinal
    \node[atom] at (5.68+\xx,\yy) {};
  \node[q label,font=\fontsize{8.5pt}{9.5pt}\selectfont] at (6.15,2.20) {$q_0,q_2$};
  \node[q label,font=\fontsize{8.5pt}{9.5pt}\selectfont] at (6.81,2.20) {$q_1,q_4$};
  \draw[galkInk,line width=.46pt] (6.04,1.91)--(6.26,1.91);
  \draw[galkInk,line width=.46pt] (6.70,1.91)--(6.92,1.91);
  % The same s_1 trap (storage row 1, column 4) remains occupied by q3.
  \node[q label,anchor=south east] at (7.19,1.20) {$q_3@s_1$};

  \draw[flow] (1.91,1.58)--(2.80,1.58);
  \node[tiny,align=center] at (2.35,2.20) {Return to\\storage};
  \node[tiny] at (2.35,1.30) {$B_0$};
  \draw[flow] (4.66,1.58)--(5.55,1.58);
  \node[tiny,align=center] at (5.10,2.20) {Place gate\\atoms};
  \node[tiny] at (5.10,1.30) {$B_1$--$B_3$};
  \end{scope}
\end{tikzpicture}
}
  \caption{双比特门层边界的联合放置实例。(a) $z_{\ell,1}$、$z_{\ell,2}$选择门位及方向；本例$\mathcal D_\ell=\{q_3\}$，$b_{\ell,1}=1$将$q_3$纳入待回迁集合，单射分配$a_\ell(q_3)=s_1$确定其存储位。(b) 精确枚举与遗传搜索共用同一编码空间。(c) 候选先将$q_3$回迁至$s_1$，更新占位后再将目标层参与原子送入所选门位；两个阶段分别检验AOD行列关系、占位顺序和非目标交点。$\mathcal E$、$\mathcal S$分别为纠缠区、存储区。}
  \label{fig:joint-ga}
\end{figure}

遗传搜索确定门位与驻留组合，匹配子问题为回迁原子分配互不重叠的存储位。对于待回迁集合$\mathcal R(\mathbf x_\ell)$中的原子$q$，从当前纠缠位、初始存储位和可见未来伙伴的邻域选取存储位，构成有界候选域：
\begin{equation}
\mathcal S_q^{\rm cand}=\operatorname{Top}
\left(\mathcal S_q^{\rm near}\cup
\mathcal S_q^{\rm home}\cup
\mathcal S_q^{\rm future}\right),
\label{eq:storage-domain}
\end{equation}
其中三个子集依次对应上述三类邻域，$\operatorname{Top}(\cdot)$按确定性代理代价仅保留有限个位置，以控制匹配规模。线性分配子问题与Murty分支给出前$K$个单射分配\cite{murty1968assignment}：
\begin{equation}
\begin{aligned}
\mathcal A_K(\mathbf x_\ell)=K\text{-Best}\{a:\mathcal R(\mathbf x_\ell)\!\rightarrow\!\mathcal S\mid{}&
a(q)\in\mathcal S_q^{\rm cand},\\[-1mm]
&a(q)\ne a(q'),~\forall q\ne q'\}.
\end{aligned}
\label{eq:kbest}
\end{equation}
每个分配$a\in\mathcal A_K$与$\mathbf x_\ell$构成完整候选。代理代价仅用于生成候选，最终次序由状态转移得到的物理代价及式~\eqref{eq:lexicographic}确定。

物理约束分为四类：占位单射、区域容量及源点释放和目标占用顺序；门位与门依赖；AOD容量、行列次序及保持关系；非目标交点。分别以$\chi_{\rm occ}$、$\chi_{\rm gate}$、$\chi_{\rm AOD}$和$\chi_{\rm int}$表示其指示函数，物理可行域为
\begin{equation}
\mathcal X_\ell^{\rm phys}=\{(\mathbf x,a)\mid
a\in\mathcal A_K(\mathbf x),
\chi_{\rm occ}\chi_{\rm gate}\chi_{\rm AOD}\chi_{\rm int}=1\}.
\label{eq:physical-domain}
\end{equation}
所选门位若被非参与原子占据，候选还需确定阻挡原子的重定位。目标层参与原子若需释放入区路径，则经可用存储位临时中转。这些辅助位置调整由候选解码确定，其阱转移、输运和时钟增量均计入物理代价。某一染色体对应的所有存储位分配均不可行时，该染色体被剔除。

\subsection{物理代价与并行重排}

物理代价取决于重排次序及轨迹的并行分组。为释放目标门位，候选先将待回迁原子及门位阻挡原子移出纠缠区，再将目标层参与原子送入所选门位。两个阶段分别形成并行重排批次，占位随已完成的移动更新；位置、资源占用和逐原子时钟均服从重排与门执行的依赖关系。对可行候选$(\mathbf x,a)$，状态转移写为
\begin{equation}
(s_{\ell+1},\Gamma_\ell)=
\mathcal T_\ell(s_\ell;\mathbf x,a),
\label{eq:state-transition}
\end{equation}
其中$\Gamma_\ell=(\mathcal B_{\ell,1},\ldots,\mathcal B_{\ell,B_\ell})$为按执行顺序排列的重排批次。$s_{\ell+1}$记录目标层Rydberg脉冲及相应单比特门完成后的状态；非参与驻留原子的空闲受激和逐原子相干损失均计入本次转换。

并行分组通过冲突图$\mathcal H_\phi=(\mathcal V_\phi,\mathcal E_\phi^{\rm AOD}\cup\mathcal E_\phi^{\rm int})$求解。顶点表示原子轨迹；违反AOD次序的轨迹之间加入$\mathcal E_\phi^{\rm AOD}$，共同激活时产生非目标交点冲突的轨迹之间加入$\mathcal E_\phi^{\rm int}$。单条轨迹若违反静态占位或交点约束，则候选不可行。其余冲突采用DSATUR着色生成初始分组\cite{brelaz1979coloring}，同色轨迹可归入同一批次，不同色轨迹分批执行。各组再依据源点释放和目标占用顺序检查，必要时继续拆分。若$\tau(\nu)$和$d(\nu)$分别为轨迹$\nu$的执行时间与输运距离，则
\begin{equation}
T_\ell^{\rm rr}=\sum_{b=1}^{B_\ell}
\max_{\nu\in\mathcal B_{\ell,b}}\tau(\nu),\qquad
D_\ell^{\rm rr}=\sum_{b=1}^{B_\ell}
\sum_{\nu\in\mathcal B_{\ell,b}}d(\nu).
\label{eq:rearrangement-metrics}
\end{equation}
批次数刻画串行重排阶段数，$T_\ell^{\rm rr}$与$D_\ell^{\rm rr}$分别刻画重排时延和总输运距离。

同一电路的单、双比特门错误在各候选间相同。边界优化比较候选相对前态的物理负对数保真度增量：
\begin{equation}
\begin{aligned}
C_\ell^{\rm cur}={}&-\Delta N_{\rm tran}\log f_{\rm tran}
-\Delta N_{\rm exc}\log f_{\rm exc}\\
&+\sum_{q\in\mathcal Q}\log\frac{1-t_q^{-}/T_2}{1-t_q^{+}/T_2}.
\end{aligned}
\label{eq:current-cost}
\end{equation}
$t_q^{-}$和$t_q^{+}$分别为状态转移前后的累计空闲时间。

\subsection{衰减多层前瞻}

衰减多层前瞻将若干层后的原子交互纳入驻留评价，从候选后状态估计未来物理代价，并赋予近层更高权重。设最大视界为$H_{\max}$，距当前边界$d$层的权重为
\begin{equation}
w_d=\alpha\rho^{d-1},\qquad 1\le d\le H_{\max}.
\label{eq:decay}
\end{equation}
记$H_{\rm cap}\le H_{\max}$为资源上限。有限视界由同时满足层距、剩余门层、资源上限与衰减阈值的未来双比特层构成，其层距集合为
\begin{equation}
\mathcal D_\ell^{\rm vis}=\{d\mid 1\le d\le H_{\rm cap},~
\mathcal G_{\ell+d}\text{存在},~\rho^{d-1}\ge\epsilon\}.
\label{eq:visible-depths}
\end{equation}
主配置取$(H_{\max},\alpha,\rho,\epsilon)=(8,0.5,0.7,0.05)$，$H_{\rm cap}$随问题规模缩减。候选后状态对下一层输运的影响及层距权重见图~\ref{fig:physical-lookahead}。

\begin{figure}[!t]
  \centering
  \resizebox{\columnwidth}{!}{% Shared vector-figure vocabulary for the Chinese IEEE manuscript.
% Every figure in this directory inputs this file, so the manuscript only
% needs the standard TikZ package and the libraries already loaded by
% paper_zh.tex.  Keep the guard: figures may be input more than once while
% editing or when a stand-alone smoke-test document is built.
\ifdefined\GALKFigureStyleLoaded\else
\def\GALKFigureStyleLoaded{1}

% Baseline-derived visual tokens: ICCAD uses a restrained blue/orange device
% palette, while ZAC uses conventional black axes and a small fixed plot set.
\definecolor{galkInk}{HTML}{231F20}
\definecolor{galkMuted}{HTML}{666666}
\definecolor{galkGrid}{HTML}{D6D6D6}
\definecolor{galkPaper}{HTML}{FFFFFF}
% Semantic colors are fixed across every figure: orange=entanglement zone,
% blue=storage zone, black=atom identity, red=invalidity, and blue dashes=future
% prediction.  Plot and search-group colours are deliberately independent of
% the two physical-zone colours.
\definecolor{galkStorage}{HTML}{0064BC}
\definecolor{galkStorageFill}{HTML}{E8F1F7}
\definecolor{galkEntangle}{HTML}{E37323}
\definecolor{galkEntangleFill}{HTML}{F9ECE3}
\definecolor{galkResident}{HTML}{485CB3}
\definecolor{galkResidentFill}{HTML}{EEF0F8}
\definecolor{galkGood}{HTML}{A4AF07}
\definecolor{galkBad}{HTML}{C84E45}
\definecolor{galkLook}{HTML}{0064BC}
\definecolor{galkData}{HTML}{315E8A}
\definecolor{galkTail}{HTML}{B44B27}
\definecolor{galkGroupOne}{HTML}{4C78A8}
\definecolor{galkGroupTwo}{HTML}{7A7A7A}
\definecolor{galkGroupThree}{HTML}{9C755F}

\tikzset{
  % Shared engineering-figure contract.  Physical snapshots must place an
  % orange entanglement band above a blue storage field; only their scale may
  % change.  Labels are aligned at the upper-left of both regions.
  galk/base/.style={
    font=\rmfamily\fontsize{9.0pt}{10.0pt}\selectfont,
    text=galkInk,
    line cap=butt,
    line join=miter,
    >=Stealth
  },
  galk/panel title/.style={
    font=\rmfamily\bfseries\fontsize{9.2pt}{10.2pt}\selectfont,
    text=galkInk,
    anchor=west
  },
  galk/label/.style={
    font=\rmfamily\fontsize{9.0pt}{10.0pt}\selectfont,
    text=galkInk
  },
  galk/small label/.style={
    font=\rmfamily\fontsize{8.7pt}{9.7pt}\selectfont,
    text=galkInk
  },
  galk/zone label/.style={
    font=\rmfamily\bfseries\fontsize{8.7pt}{9.7pt}\selectfont,
    text=galkInk,
    anchor=north west,
    inner sep=.5pt
  },
  galk/entanglement zone/.style={
    draw=galkEntangle,
    fill=galkEntangleFill,
    line width=.40pt
  },
  galk/storage zone/.style={
    draw=galkStorage,
    fill=galkStorageFill,
    line width=.40pt
  },
  galk/site/.style={
    circle,
    draw=galkInk!68,
    fill=white,
    line width=.36pt,
    minimum size=1.35mm,
    inner sep=0pt
  },
  galk/occupied atom/.style={
    circle,
    draw=galkInk,
    fill=galkInk,
    line width=.36pt,
    minimum size=1.55mm,
    inner sep=0pt
  },
  galk/target site/.style={
    circle,
    draw=galkInk!60,
    fill=galkInk!10,
    line width=.34pt,
    minimum size=1.55mm,
    inner sep=0pt
  },
  galk/q label/.style={
    font=\rmfamily\itshape\fontsize{8.5pt}{9.5pt}\selectfont,
    text=galkInk,
    inner sep=.25pt
  },
  galk/rydberg pair/.style={
    draw=galkInk!62,
    densely dashed,
    line width=.36pt
  },
  galk/accepted path/.style={
    -{Stealth[length=.95mm,width=.62mm]},
    draw=galkInk,
    line width=.48pt
  },
  galk/alternative path/.style={
    -{Stealth[length=.95mm,width=.62mm]},
    draw=galkMuted,
    densely dashed,
    line width=.48pt
  },
  galk/future path/.style={
    -{Stealth[length=.95mm,width=.62mm]},
    draw=galkLook,
    densely dashed,
    line width=.48pt
  },
  galk/invalid path/.style={
    -{Stealth[length=.95mm,width=.62mm]},
    draw=galkBad,
    densely dashed,
    line width=.50pt
  }
}
\fi

\begin{tikzpicture}[galk/base,x=1cm,y=1cm,
  panel title/.style={galk/panel title},
  small label/.style={galk/small label},
  zone label/.style={galk/zone label},
  site/.style={galk/site},
  occupied/.style={galk/occupied atom},
  target/.style={galk/target site},
  route/.style={galk/accepted path},
  graph vertex/.style={circle,draw=galkInk,line width=.42pt,
                       minimum size=2.25mm,inner sep=0pt}
]
  \path[use as bounding box] (0,-.22) rectangle (8.50,8.22);

  % (a) Physical trajectories, their conflict graph, and the resulting groups.
  \node[panel title] at (.08,7.98) {(a) Compatibility at $\pi_\ell\!\to\!\pi_{\ell+1}$};
  \node[small label] at (1.50,7.62) {trajectories};
  \node[small label,anchor=east] at (.53,7.50) {$\pi_\ell$};
  \node[small label,anchor=west] at (2.54,7.50) {$\pi_{\ell+1}$};

  \coordinate (s1) at (.58,7.20);
  \coordinate (s2) at (.58,6.82);
  \coordinate (s3) at (.58,6.42);
  \coordinate (s4) at (.58,6.04);
  \coordinate (t1) at (2.48,6.82);
  \coordinate (t2) at (2.48,7.20);
  \coordinate (t3) at (2.48,6.04);
  \coordinate (t4) at (2.48,6.42);
  \foreach \i in {1,2,3,4} {
    \node[occupied] at (s\i) {};
    \node[target] at (t\i) {};
    \node[small label,anchor=east] at ([xshift=-1.4pt]s\i) {$q_{\i}$};
    \node[small label,anchor=west] at ([xshift=1.4pt]t\i) {$q_{\i}$};
  }
  \draw[route] (s1)--(t1);
  \draw[route] (s2)--(t2);
  \draw[route] (s3)--(t3);
  \draw[route] (s4)--(t4);

  \draw[-{Stealth[length=.85mm,width=.55mm]},galkInk,line width=.42pt]
    (2.95,6.61)--(3.35,6.61);
  \node[small label] at (4.25,7.62) {conflict graph};

  \node[graph vertex,fill=galkGroupTwo!18] (v1) at (3.68,7.18) {};
  \node[graph vertex,fill=galkGroupOne!18] (v2) at (4.15,6.43) {};
  \node[graph vertex,fill=galkGroupTwo!18] (v3) at (4.83,6.94) {};
  \node[graph vertex,fill=galkGroupOne!18] (v4) at (4.89,6.10) {};
  \node[small label,anchor=east] at ([xshift=-1.5pt]v1.west) {$q_1$};
  \node[small label,anchor=east] at ([xshift=-1.5pt]v2.west) {$q_2$};
  \node[small label,anchor=west] at ([xshift=1.5pt]v3.east) {$q_3$};
  \node[small label,anchor=west] at ([xshift=1.5pt]v4.east) {$q_4$};
  \draw[galkBad,line width=.46pt] (v1)--(v2);
  \draw[galkBad,line width=.46pt] (v3)--(v4);
  \node[small label] at (4.25,5.73)
    {$\{q_1,q_2\},\ \{q_3,q_4\}$};

  \draw[-{Stealth[length=.85mm,width=.55mm]},galkInk,line width=.42pt]
    (5.34,6.61)--(5.77,6.61);
  \node[small label] at (6.88,7.62) {DSATUR coloring};
  \node[small label,anchor=west] at (5.92,7.20)
    {$\mathcal B_1=\{q_2,q_4\}$};
  \node[small label,anchor=west] at (5.92,6.67)
    {$\mathcal B_2=\{q_1,q_3\}$};
  \draw[galkGroupOne,line width=1.05pt] (5.92,7.03)--(6.32,7.03);
  \draw[galkGroupTwo,line width=1.05pt] (5.92,6.50)--(6.32,6.50);
  \node[small label,anchor=west,align=left] at (5.92,6.10)
    {$q_0$: stationary\\no trajectory};

  \draw[galkInk!32,line width=.35pt] (.08,5.47)--(8.20,5.47);

  % (b) Equal immediate costs may lead to different finite-horizon costs.
  \node[panel title] at (.08,5.31) {(b) Dependence on the post-layer state};
  \node[small label] at (4.17,4.96)
    {$C_{\ell}^{\rm cur}(A)=C_{\ell}^{\rm cur}(B)$};

  \foreach \xzero in {.42,4.65} {
    \path[galk/entanglement zone]
      (\xzero,3.70) rectangle (\xzero+3.22,4.64);
    \path[galk/storage zone]
      (\xzero,2.08) rectangle (\xzero+3.22,3.54);
    \node[zone label,fill=galkEntangleFill]
      at (\xzero+.08,4.58) {$\mathcal E$};
    \node[zone label,fill=galkStorageFill]
      at (\xzero+.08,3.48) {$\mathcal S$};
    \foreach \dx in {.38,.98,1.58,2.18,2.78} {
      \foreach \yy in {2.38,2.80,3.22} {
        \node[site] at (\xzero+\dx,\yy) {};
      }
    }
    \foreach \px in {.72,2.18} {
      \node[site] at (\xzero+\px,3.95) {};
      \node[site] at (\xzero+\px+.26,3.95) {};
      \draw[galk/rydberg pair]
        (\xzero+\px+.13,3.95) ellipse (.21 and .13);
    }
  }
  \node[small label] at (1.10,4.94)
    {$A:\ s_{\ell+1}$};
  \node[small label] at (7.18,4.94)
    {$B:\ s_{\ell+1}$};

  % Candidate A: q and r occupy the current gate pair.  The next layer moves
  % q and p to the second pair after r returns to storage.
  \node[occupied] (ar) at (1.14,3.95) {};
  \node[occupied] (aq) at (1.40,3.95) {};
  \node[target] (aqt) at (2.60,3.95) {};
  \node[target] (apt) at (2.86,3.95) {};
  \node[occupied] (ap) at (3.20,2.38) {};
  \node[target] (art) at (2.00,3.22) {};
  \node[galk/q label,anchor=north east] at ([xshift=-2.8pt,yshift=-1.4pt]ar) {$r$};
  \node[galk/q label,anchor=north west] at ([xshift=2.8pt,yshift=-1.4pt]aq) {$q$};
  \node[galk/q label,anchor=west] at ([xshift=3.2pt]ap) {$p$};
  \node[small label] at (1.27,4.18) {$\mathrm{CZ}(q,r)$};
  \node[small label] at (2.73,4.18) {$\mathrm{CZ}(q,p)$};
  \draw[galk/alternative path]
    (ar) .. controls (1.45,3.72) and (1.76,3.48) .. (art);
  \draw[galk/future path] (aq)--(aqt);
  \draw[galk/future path]
    (ap) .. controls (3.25,3.12) and (3.10,3.57) ..
    node[pos=.44,left,small label] {$d_A$} (apt);

  % Candidate B has the same current gate pair and target pair, but places p
  % nearer the target used by CZ(q,p).
  \node[occupied] (br) at (5.37,3.95) {};
  \node[occupied] (bq) at (5.63,3.95) {};
  \node[target] (bqt) at (6.83,3.95) {};
  \node[target] (bpt) at (7.09,3.95) {};
  \node[occupied] (bp) at (6.83,3.22) {};
  \node[target] (brt) at (6.23,3.22) {};
  \node[galk/q label,anchor=north east] at ([xshift=-2.8pt,yshift=-1.4pt]br) {$r$};
  \node[galk/q label,anchor=north west] at ([xshift=2.8pt,yshift=-1.4pt]bq) {$q$};
  \node[galk/q label,anchor=west] at ([xshift=3.2pt]bp) {$p$};
  \node[small label] at (5.50,4.18) {$\mathrm{CZ}(q,r)$};
  \node[small label] at (6.96,4.18) {$\mathrm{CZ}(q,p)$};
  \draw[galk/alternative path]
    (br) .. controls (5.68,3.72) and (5.99,3.48) .. (brt);
  \draw[galk/future path] (bq)--(bqt);
  \draw[galk/future path]
    (bp)--node[pos=.70,right,small label,fill=white,inner sep=.2pt] {$d_B$} (bpt);

  \node[small label] at (4.17,1.82)
    {$d_B<d_A\ \Rightarrow\ C_{\ell,1}^{\rm gate}(B)<C_{\ell,1}^{\rm gate}(A)$};

  \draw[galkInk!32,line width=.35pt] (.08,1.55)--(8.20,1.55);

  % (c) Conventional axis representation of the geometric decay.
  \node[panel title] at (.08,1.31) {(c) Decayed layer weights};
  \draw[-{Stealth[length=.8mm,width=.52mm]},galkInk,line width=.42pt]
    (.72,.27)--(5.14,.27) node[right,small label] {$d$};
  \draw[-{Stealth[length=.8mm,width=.52mm]},galkInk,line width=.42pt]
    (.72,.27)--(.72,1.05);
  \node[small label,rotate=90] at (.49,.69) {$w_d$};
  \draw[galkData,line width=.55pt]
    (1.23,1.02)--(2.08,.82)--(2.93,.66)--(3.78,.53)--(4.63,.43);
  \foreach \x/\y/\lab in {1.23/1.02/1,2.08/.82/2,2.93/.66/3,
                            3.78/.53/\cdots,4.63/.43/d_\ell^*} {
    \node[circle,draw=galkData,fill=white,line width=.45pt,
          minimum size=1.35mm,inner sep=0pt] at (\x,\y) {};
    \draw[galkInk,line width=.35pt] (\x,.23)--(\x,.31);
    \node[small label,anchor=north] at (\x,.20) {$\lab$};
  }
  \node[small label,anchor=west] at (5.38,.72)
    {$w_d=\alpha\rho^{d-1}$};
\end{tikzpicture}
}
  \caption{物理可行性与有限视界评价实例。(a) 四条轨迹的冲突图着色产生$\mathcal B_1=\{q_2,q_4\}$和$\mathcal B_2=\{q_1,q_3\}$；实心、空心圆分别为源点和目标点，静止的$q_0$不入图。(b) 当前代价相同的候选$A$、$B$导致不同后状态；在其余物理增量相同的图示条件下，$d_B<d_A$使候选$B$的下一层输运分量较低。蓝色和灰色虚线分别表示下一层与回迁轨迹，$\mathcal E$、$\mathcal S$分别为纠缠区、存储区。(c) 层距权重按$w_d=\alpha\rho^{d-1}$在实际可见层集合$\mathcal D_\ell^{\rm vis}$内衰减。}
  \label{fig:physical-lookahead}
\end{figure}

未来代价从各候选的后状态$s'=s_{\ell+1}$分别估计。令$s_{\ell,0}=s'$，确定性状态转移$\mathcal F$逐层给出未来状态及增量代价：
\begin{equation}
(s_{\ell,d},\widehat C_{\ell,d})=
\mathcal F(s_{\ell,d-1},\mathcal G_{\ell+d}),
\qquad d\in\mathcal D_\ell^{\rm vis},
\label{eq:forecast-transition}
\end{equation}
其中，$\mathcal F$为本预测层各门选择尚未分配的合法纠缠位对，并将阻挡目标端点或轨迹的非参与原子迁往存储区。第$d$层的未加权物理增量为
\begin{equation}
\widehat C_{\ell,d}=C_{\ell,d}^{\rm res}
+C_{\ell,d}^{\rm gate}+C_{\ell,d}^{\rm block},
\label{eq:forecast-components}
\end{equation}
其中，$C_{\ell,d}^{\rm res}$为非参与驻留原子的空闲受激和相干损失；$C_{\ell,d}^{\rm gate}$为门参与原子入区及临时中转产生的物理损失；$C_{\ell,d}^{\rm block}$为阻挡原子回迁产生的物理损失。后两项均包含阱转移、输运与相干损失。

未来门$(u,v)$的门位由两参与原子的移动距离与端点阻挡原子的最近存储距离之和确定，同分时按门位及方向编号选择。联合入区不可行时，依次考察有限中转与逐原子入区方案。若均不可行，则该候选为预测不可行候选；其当前边界评价仍保留，用于无预测可行候选时的回退选择。

若$\mathcal D_\ell^{\rm vis}$非空，令$d_\ell^*=\max\mathcal D_\ell^{\rm vis}$；窗口末端以清理势$\Phi(s_{\ell,d_\ell^*})$计入尚未完成的驻留代价：
\begin{equation}
\widehat C_\ell^{\rm future}=
\sum_{d\in\mathcal D_\ell^{\rm vis}}\alpha\rho^{d-1}\widehat C_{\ell,d}
+\alpha\rho^{d_\ell^*-1}\Phi(s_{\ell,d_\ell^*}).
\label{eq:lookahead-cost}
\end{equation}
当$H_{\max}=0$时，$\mathcal D_\ell^{\rm vis}=\varnothing$且$\widehat C_\ell^{\rm future}=0$。

终端清理势$\Phi$估计窗口结束后尚需承担的回迁损失：对仍留在纠缠区、且未参与窗口末层的原子构造可行回迁，取其负对数保真度增量；可见层集合为空时取零。

\subsection{候选排序与边界选择}

未来收益可能伴随当前代价增加，候选选择需控制两者的交换幅度。求解器以当前代价较低的候选为锚点，限定其他候选的当前代价增量，再按有限视界目标排序：
\begin{equation}
\mathbf J_\ell=(C_\ell^{\rm cur}+\widehat C_\ell^{\rm future},B_\ell,
T_\ell^{\rm rr},D_\ell^{\rm rr},\kappa_\ell)
\label{eq:lexicographic}
\end{equation}
词典序先比较首项，同分时再依次比较后续项。$\kappa_\ell=(\overline{\mathbf x}_\ell,\overline a_\ell)$为染色体与存储位分配的规范键，用于确定完全同分候选的次序。总负对数代价、时延和距离分别按$10^{-12}$、$10^{-6}\,\mu\mathrm{s}$和$10^{-6}\,\mu\mathrm{m}$量化，批次数采用整数比较。

锚点的构造包含门位投影与驻留保护。门位投影在主搜索后固定已评价候选的驻留编码段，逐坐标调整门位基因以降低当前代价。采用的驻留编码包括当前最优、全驻留以及由物理往返比较得到的回迁或驻留建议。规模受限时扫描一轮有界门位选项，其余情形至多扫描两轮；每个试探候选均重新归一化并检查物理可行性。投影结果与采用相同驻留编码的已评价候选构成邻域$\mathcal N_\ell^{\rm guard}$，新增评价数记为$E_{\rm guard}$。

驻留保护比较原子在可见复用层之前的驻留损失与往返存储区损失。前者计入空闲受激与相干损失，后者计入可行路径上的阱转移、输运与相干损失。驻留代价较低且满足强制回迁和容量约束的原子组成保护集合$\mathcal P_\ell$，候选选择优先保留这些原子的驻留；$H_{\max}=0$时该集合为空。候选$u=(\mathbf x,a)$回迁保护集合内原子的数量为
\begin{equation}
v_\ell(u)=\sum_{q_i\in\mathcal P_\ell}
\mathbf 1\!\left[b_{\ell,i}(u)=1\right].
\label{eq:stay-violation}
\end{equation}
记$\mathcal X_\ell^{\rm eval}\subseteq\mathcal X_\ell^{\rm phys}$为已完成评价的物理可行候选集合。评价包括当前状态转换，以及启用前瞻时的未来层预测。求解器按$v_\ell$递增检查候选组，取首个含预测可行候选的组，并按当前物理键$\mathbf J_\ell^{\rm cur}=(C_\ell^{\rm cur},B_\ell,T_\ell^{\rm rr},D_\ell^{\rm rr},\overline{\mathbf x}_\ell)$确定锚点$u^0$。令$L_\ell(u)=C_\ell^{\rm cur}(u)+\widehat C_\ell^{\rm future}(u)$，并令$\Delta C_{\rm cur}(u)=\max\{0,C_\ell^{\rm cur}(u)-C_\ell^{\rm cur}(u^0)\}$、$\Delta S_{\rm future}(u)=\max\{0,\widehat C_\ell^{\rm future}(u^0)-\widehat C_\ell^{\rm future}(u)\}$。准入集合为
\begin{equation}
\begin{aligned}
\mathcal X_\ell^{\rm adm}=\{u\in\mathcal X_\ell^{\rm eval}:{}&
v_\ell(u)=v_\ell(u^0),\\[-1mm]
&\Delta C_{\rm cur}(u)\le\eta\Delta S_{\rm future}(u),\\[-1mm]
&L_\ell(u)\le L_\ell(u^0)\}.
\end{aligned}
\label{eq:admission-domain}
\end{equation}
其中$\eta=0.25$，即当前代价增量至多为预计未来节省的四分之一。准入集合是物理可行域的子集。
若不存在通过有限视界可行性检查的候选，则在当前边界可行候选中按$\mathbf J_\ell^{\rm cur}$选择。存在预测可行候选时，在准入集合上进行词典序选择：
\begin{equation}
(\mathbf x_\ell^*,a_\ell^*)=
\mathop{\mathrm{lexarg\,min}}_{(\mathbf x,a)\in\mathcal X_\ell^{\rm adm}}
\mathbf J_\ell(\mathbf x,a).
\label{eq:boundary-objective}
\end{equation}
该规则优先最小化当前与未来代价之和，同分时依次比较批次数、重排时延、输运距离和规范编码。

\subsection{精确枚举与遗传搜索}

小规模边界采用完整枚举，大规模边界采用遗传搜索。边界编码集合为$\mathcal Z_\ell=\prod_j\{0,\ldots,|\Omega_{\ell,j}|-1\}\times\{0,1\}^{|\mathcal D_\ell|}$，编码数为$D_\ell=|\mathcal Z_\ell|$。当$D_\ell\le\min(D_{\rm enum},E_{\max})$时枚举全部编码，否则采用有界遗传搜索\cite{holland1975adaptation}；$D_{\rm enum}$和$E_{\max}$分别为枚举阈值与不同染色体评价预算。两种方式均通过式~\eqref{eq:kbest}分配回迁存储位，并采用相同的物理评价和准入规则。

遗传搜索的确定性种子集$\mathcal S_\ell^0$包含全驻留编码、全回迁编码、物理贪心编码和回迁建议编码。若LRU缓存中存在与当前完整边界状态键相同的染色体，则优先使用。门位段与驻留段分别进行单点交叉；未交叉的子代采用随机单基因变异、高输运代价门位重选、门位--驻留耦合变异或回迁成本引导的驻留位翻转。候选经约束归一化后按染色体键去重，已评价染色体的目标值由缓存复用。

搜索在达到$E_{\max}$、完成$I$代或最优值连续$p_{\rm stop}$代未改善时终止；预算有余时，对当前最优解执行有限轮局部邻域下降。随机种子固定以保证可复现性。记$\mathcal U_\ell=(s_\ell,\mathcal G_\ell)$为当前边界实例，$\Theta$汇集区域与AOD参数、物理模型常数、候选域上限和求解预算。完整求解过程见算法~\ref{alg:solver}。

\begin{algorithm}[t]
\caption{层边界联合放置}
\label{alg:solver}
\scriptsize
\begin{algorithmic}[1]
\Procedure{Evaluate}{$\mathcal U_\ell,\mathbf x,\mathcal D_\ell^{\rm vis},\Theta$}
  \State $\mathcal A\gets\Call{StorageAssignment}{\mathcal U_\ell,\mathbf x}$；$\mathcal R_{\mathbf x}\gets\varnothing$
  \For{每个存储位分配$a\in\mathcal A$}
    \State 导出阻挡原子重定位、临时中转和原子轨迹
    \If{占位、容量、AOD或交点约束不满足}
      \State \textbf{continue}
    \EndIf
    \State 对冲突图着色，并确定可执行的重排分组及后状态$s'$
    \State 计算当前代价$C^{\rm cur}$与重排指标$B,T,D$
    \State $(\widehat C^{\rm future},\zeta)\gets\Call{FutureCost}{s',\mathcal D_\ell^{\rm vis},\Theta}$
    \State 将当前物理记录及前瞻可行标志$\zeta$加入$\mathcal R_{\mathbf x}$
  \EndFor
  \State \textbf{输出} $\mathcal R_{\mathbf x}$
\EndProcedure
\Statex
\Procedure{OptimizeBoundary}{$\mathcal U_\ell,H_{\max},\Theta$}
  \State 计算编码上界$D_\ell$
  \State $\mathcal D_\ell^{\rm vis}\gets\Call{VisibleDepths}{\mathcal U_\ell,H_{\max},\Theta}$
  \State $\mathcal R\gets\varnothing$
  \If{$D_\ell\le\min(D_{\rm enum},E_{\max})$}
    \For{每个$\mathbf x\in\mathcal Z_\ell$}
      \State $\mathcal R\gets\mathcal R\cup\Call{Evaluate}{\mathcal U_\ell,\mathbf x,\mathcal D_\ell^{\rm vis},\Theta}$
    \EndFor
  \Else
    \State $\mathcal C\gets\Call{NormalizeDedup}{\mathcal S_\ell^0,\text{缓存种子、随机种子}}$
    \For{$g=1,\ldots,I$}
      \State 生成并评价交叉或变异子代，将可行记录加入$\mathcal R$
      \State 保留目标值最优的$P$个规范染色体
      \If{评价预算耗尽或连续$p_{\rm stop}$代未改善}
        \State \textbf{break}
      \EndIf
    \EndFor
  \EndIf
  \State 对$\mathcal R$中的当前最优解进行有限轮局部精修
  \If{$\mathcal R$中不存在$\zeta=1$的记录}
    \State \textbf{输出} 当前物理键最小的记录，并置$\widehat C^{\rm future}=0$
  \EndIf
  \State 合并主搜索与门位--驻留保护邻域中的已评价记录，构造$\mathcal X_\ell^{\rm adm}$
  \State \textbf{输出}$\mathcal X_\ell^{\rm adm}$中$\mathbf J_\ell$最小的候选
\EndProcedure
\end{algorithmic}
\end{algorithm}

\subsection{可行性与求解性质}

候选生成、投影与准入选择共用$\mathcal X_\ell^{\rm phys}$的约束。每次接受的状态转换均满足物理约束及重排批次的顺序依赖，占位单射、区域容量、AOD次序、门依赖和逐原子时钟的一致性在各层转换中保持。

候选投影与准入阶段未激活且评价预算覆盖全部有界候选时，完整枚举返回$\mathcal X_\ell^{\rm phys}$上的词典序最优解；该阶段激活时返回$\mathcal X_\ell^{\rm adm}$中的最优候选。遗传搜索在固定预算内终止，不保证全局最优。

\subsection{复杂度与实现}

实现采用C++17，并预计算坐标距离、轨迹兼容关系和冲突边；候选解码、存储位分配、冲突图着色及未来状态均使用确定性缓存。令$N_{\rm bnd}$表示相邻双比特层的边界数；表~\ref{tab:solver-bounds}给出按$N_{\rm bnd}$分段的求解预算，用于限制长电路的候选评价、存储位匹配和预测层数。

\begin{table}[!t]
\caption{GA-LK的默认及规模自适应求解上限}
\label{tab:solver-bounds}
\centering
\PaperTableSetup
\begin{tabular*}{\columnwidth}{@{\extracolsep{\fill}}lccc@{}}
\toprule
配置 & 枚举上限 & $(E,P,I,K)$ & $H_{\rm cap}$\\
\midrule
默认配置 & 512 & $(576,8,6,4)$ & 8\\
$N_{\rm bnd}>512$ & 64 & $(192,\le6,\le4,\le2)$ & $\le4$\\
\shortstack[l]{$N_{\rm bnd}\ge5000$\\$n\le16$} & 16 & $(32,\le4,\le1,1)$ & --\\
$N_{\rm bnd}\ge5000$ & -- & -- & $\le2$\\
\bottomrule
\end{tabular*}
\PaperTableNote{注：$E$、$P$、$I$和$K$分别为候选评价预算、种群规模、迭代数和单射分配数；“--”表示该行不新增对应项约束，重叠条件同时生效并取更严格上限。}
\end{table}

其中$n=|\mathcal Q|$；各原子的存储位候选域至多含6个位置。冲突图精确改进至多处理24条轨迹，分支限界节点上限为$2\times10^5$。存储位单射分配由线性分配子问题与Murty分支生成，轨迹两两兼容关系的构造对轨迹数呈二次增长。

通常，$\Omega_{\ell,j}$包含通过静态占位与单轨迹交点检查的全部纠缠位对。对于深度较大且单层门数较少的电路，首轮求解使用按距离排序的有界门位域；未产生可行候选时恢复完整域。表~\ref{tab:solver-bounds}与代码仓库给出规模阈值和完整预算。缓存以完整边界状态和规范编码为键；最终重排指标由生成的ZAIR指令序列计算。

设主搜索和候选投影阶段分别评价$E$和$E_{\rm guard}$个候选，每个候选考虑至多$K$个存储位分配，并展开$|\mathcal D_\ell^{\rm vis}|$个未来层，则时间复杂度近似为
\begin{equation}
O\!\left((E+E_{\rm guard})K
\left(C_{\rm feas}+|\mathcal D_\ell^{\rm vis}|C_{\rm pred}\right)\right),
\label{eq:complexity}
\end{equation}
其中$C_{\rm feas}$和$C_{\rm pred}$分别表示物理可行性评价与单层未来状态估计的代价；宽层和长深度实例对$E_{\rm guard}$采用有界门位邻域。生成的ZAIR指令序列需满足占位、AOD和非目标交点约束。
